%% ============================================================
%%  OLIBOT — Presentación para Dirección de Centro Escolar
%%  Objetivo: obtener autorización para evaluación con alumnos
%%  Compilar: pdflatex presentacion_colegio.tex  (x2)
%% ============================================================
\documentclass[aspectratio=169, 11pt]{beamer}

\usepackage[spanish]{babel}
\usepackage[utf8]{inputenc}
\usepackage[T1]{fontenc}
\usepackage{lmodern}
\usepackage{tikz}
\usepackage{booktabs}
\usepackage{xcolor}
\usepackage{graphicx}
\usepackage{microtype}
\usepackage{fontawesome5}
\usepackage{array}
\usepackage{colortbl}

\usetikzlibrary{arrows.meta, positioning, shapes.geometric, fit, backgrounds, shadows}

%% ── Paleta amigable ──────────────────────────────────────────
\definecolor{verde}   {RGB}{34,  139,  86}
\definecolor{azul}    {RGB}{30,   90, 160}
\definecolor{naranja} {RGB}{230, 120,  30}
\definecolor{crema}   {RGB}{252, 248, 240}
\definecolor{grisclaro}{RGB}{245,245,245}
\definecolor{verdeOsc}{RGB}{20, 100, 60}

%% ── Tema ─────────────────────────────────────────────────────
\usetheme{Boadilla}
\usecolortheme[named=verde]{structure}

\setbeamercolor{title}             {fg=white,  bg=verde}
\setbeamercolor{frametitle}        {fg=white,  bg=verde}
\setbeamercolor{block title}       {fg=white,  bg=azul}
\setbeamercolor{block body}        {bg=azul!7}
\setbeamercolor{block title alerted}{fg=white, bg=naranja}
\setbeamercolor{block body alerted} {bg=naranja!8}
\setbeamercolor{block title example}{fg=white, bg=verdeOsc}
\setbeamercolor{block body example} {bg=verde!8}
\setbeamercolor{itemize item}      {fg=verde}
\setbeamercolor{itemize subitem}   {fg=azul}
\setbeamercolor{enumerate item}    {fg=verde}
\setbeamercolor{footline}          {fg=white,  bg=verde!80}
\setbeamercolor{palette primary}   {fg=white,  bg=verde}
\setbeamercolor{palette secondary} {fg=white,  bg=verde!80}
\setbeamercolor{palette tertiary}  {fg=white,  bg=verde!60}

\setbeamertemplate{navigation symbols}{}
\setbeamertemplate{itemize items}[circle]
\setbeamertemplate{itemize subitem}[triangle]
\setbeamerfont{frametitle}{size=\large, series=\bfseries}

%% Barra de pie personalizada
\setbeamertemplate{footline}{%
  \leavevmode%
  \hbox{%
    \begin{beamercolorbox}[wd=.40\paperwidth, ht=2.4ex, dp=1ex, leftskip=6pt]{palette primary}%
      \usebeamerfont{author in head/foot}\insertshortauthor
    \end{beamercolorbox}%
    \begin{beamercolorbox}[wd=.45\paperwidth, ht=2.4ex, dp=1ex, center]{palette secondary}%
      \usebeamerfont{title in head/foot}\insertshorttitle
    \end{beamercolorbox}%
    \begin{beamercolorbox}[wd=.15\paperwidth, ht=2.4ex, dp=1ex, rightskip=6pt, right]{palette tertiary}%
      \usebeamerfont{date in head/foot}\insertframenumber\,/\,\inserttotalframenumber
    \end{beamercolorbox}%
  }%
}

%% ── Macros ───────────────────────────────────────────────────
\newcommand{\icono}[1]{\textcolor{verde}{\faIcon{#1}}}
\newcommand{\destacado}[1]{\textbf{\textcolor{azul}{#1}}}
\newcommand{\tick}{\textcolor{verde}{\faIcon{check-circle}}}
\newcommand{\punto}[1]{\medskip\noindent\textcolor{verde}{\faIcon{circle}}\;#1}
\newcommand{\titulo}{\bfseries\color{azul}}

%% ── Metadatos ────────────────────────────────────────────────
\title[OLIBOT — Asistente Digital Infantil]{%
  \textbf{OLIBOT}\\[6pt]
  \large Un asistente digital que apoya\\[2pt]
  el aprendizaje de la lecto\-escritura\\[2pt]
  en Educación Infantil}
\author[María Emilia Núñez Guerrero]{%
  María Emilia Núñez Guerrero\\[2pt]
  \small Máster en Inteligencia Artificial — Universidad Carlos III de Madrid}
\institute{}
\date{Junio 2026}

%% ============================================================
\begin{document}
%% ============================================================

%% ─── 1. PORTADA ─────────────────────────────────────────────
\begin{frame}
  \begin{tikzpicture}[remember picture, overlay]
    \fill[verde!15] (current page.south west) rectangle (current page.north east);
    \fill[verde]    (current page.north west) rectangle
                    ([yshift=-1.8cm]current page.north east);
    \node[white, font=\bfseries\huge, anchor=west]
      at ([xshift=0.8cm, yshift=-0.9cm]current page.north west)
      {OLIBOT};
    \node[white, font=\normalsize, anchor=east]
      at ([xshift=-0.8cm, yshift=-0.9cm]current page.north east)
      {Junio 2026};
  \end{tikzpicture}

  \vspace{1.6cm}
  \begin{center}
    {\Huge\bfseries\color{verde} OLIBOT}\\[6pt]
    {\large\color{azul} Un asistente digital para aprender\\
    las primeras letras y números}\\[18pt]
    \begin{tikzpicture}
      \node[fill=verde, rounded corners=8pt, inner sep=10pt, text=white,
            font=\bfseries\large] {
        Propuesta de evaluación en el aula};
    \end{tikzpicture}\\[14pt]
    {\normalsize María Emilia Núñez Guerrero}\\
    {\small TFM — Universidad Carlos III de Madrid}
  \end{center}
\end{frame}

%% ─── 2. ÍNDICE ───────────────────────────────────────────────
\begin{frame}{¿De qué vamos a hablar?}
  \tableofcontents
\end{frame}

%% ============================================================
\section{El contexto educativo actual}
%% ============================================================

%% ─── 3. NORMATIVA CM ────────────────────────────────────────
\begin{frame}{El currículo de Infantil en la Comunidad de Madrid}
  \begin{columns}[T]
    \begin{column}{.54\textwidth}
      \begin{block}{\icono{book-open}\; Decreto 36/2022 de la CM}
        \small
        El currículo de 2.º ciclo de Educación Infantil (3--6 años)
        establece que la \destacado{iniciación a la lectoescritura}
        es un eje central del área\\
        \textit{«Comunicación y Representación de la Realidad»}:
        \begin{itemize}
          \item Grafomotricidad y trazado de letras desde los 3 años.
          \item Conciencia fonológica y asociación grafema–fonema.
          \item Aproximación al número y sus representaciones.
          \item Enfoque \textbf{sistemático y progresivo} por edades.
        \end{itemize}
      \end{block}
    \end{column}
    \begin{column}{.42\textwidth}
      \begin{alertblock}{\icono{exclamation-triangle}\; El reto del aula}
        \small
        \begin{itemize}
          \item Ratios de \textbf{25 alumnos/docente}.
          \item Ritmos de aprendizaje \textbf{muy distintos}.
          \item Escaso tiempo para atención \textbf{individualizada}.
          \item El juego y la motivación son \textbf{imprescindibles} a estas edades.
        \end{itemize}
      \end{alertblock}
      \vspace{4pt}
      \begin{exampleblock}{\icono{lightbulb}\; La pregunta}
        \small\centering
        ¿Puede la tecnología \textbf{complementar} al docente respetando el currículo?
      \end{exampleblock}
    \end{column}
  \end{columns}
\end{frame}

%% ─── 4. DIGITAL HOY ─────────────────────────────────────────
\begin{frame}{Los niños de hoy y la tecnología}
  \begin{columns}[T]
    \begin{column}{.48\textwidth}
      \begin{block}{\icono{child}\; Nativos digitales}
        \footnotesize
        Los niños de 3--5 años interactúan con dispositivos
        táctiles y de voz a diario.\\[2pt]
        Incorporar la tecnología de forma \destacado{controlada y pedagógica}:
        \begin{itemize}
          \item Reduce la brecha con su entorno familiar.
          \item Aumenta la motivación intrínseca.
          \item Permite practicar a su propio ritmo.
        \end{itemize}
      \end{block}
      \begin{block}{\icono{shield-alt}\; Diferencia clave}
        \footnotesize
        OLIBOT \textbf{no es entretenimiento}: es una herramienta educativa con objetivos curriculares, sin publicidad ni datos comerciales.
      \end{block}
    \end{column}
    \begin{column}{.48\textwidth}
      \begin{tikzpicture}[font=\footnotesize]
        % Comparativa Apps vs OLIBOT
        \node[draw=naranja, fill=naranja!10, rounded corners=5pt,
              text width=4.0cm, inner sep=6pt, align=left] (apps)
          {{\color{naranja}\bfseries Apps típicas}\\[4pt]
           \textcolor{gray!80}{\small\faIcon{times-circle}} Contenido fijo\\
           \textcolor{gray!80}{\small\faIcon{times-circle}} No se adapta al niño\\
           \textcolor{gray!80}{\small\faIcon{times-circle}} Sin alineación curricular\\
           \textcolor{gray!80}{\small\faIcon{times-circle}} Refuerzo instantáneo adictivo};
        \node[draw=verde, fill=verde!10, rounded corners=5pt,
              text width=4.0cm, inner sep=6pt, align=left,
              below=0.35cm of apps] (olibot)
          {{\color{verde}\bfseries OLIBOT}\\[4pt]
           \textcolor{verde}{\small\faIcon{check-circle}} Adaptativo por edad y nivel\\
           \textcolor{verde}{\small\faIcon{check-circle}} Sigue el currículo CM\\
           \textcolor{verde}{\small\faIcon{check-circle}} Interacción por voz\\
           \textcolor{verde}{\small\faIcon{check-circle}} Sin publicidad ni datos personales};
      \end{tikzpicture}
    \end{column}
  \end{columns}
\end{frame}

%% ============================================================
\section{¿Qué es OLIBOT?}
%% ============================================================

%% ─── 5. QUÉ ES ───────────────────────────────────────────────
\begin{frame}{OLIBOT en pocas palabras}
  \begin{columns}[T]
    \begin{column}{.50\textwidth}
      \begin{exampleblock}{\icono{robot}\; Qué es}
        \small
        Un \destacado{asistente educativo digital} que habla con el niño por voz, propone actividades de lectoescritura y se adapta a su ritmo.\\[3pt]
        Funciona en cualquier \textbf{tablet} (sin instalación).
      \end{exampleblock}
      \vspace{3pt}
      \begin{block}{\icono{users}\; Para quién}
        \footnotesize
        \begin{itemize}
          \item \textbf{3 años}: grafomotricidad y vocales
          \item \textbf{4 años}: vocales + primeras consonantes
          \item \textbf{5 años}: consonantes, sílabas y palabras
        \end{itemize}
        El nivel se selecciona al inicio; la dificultad se ajusta automáticamente.
      \end{block}
    \end{column}
    \begin{column}{.46\textwidth}
      \begin{tikzpicture}[font=\footnotesize,
        fase/.style={draw=azul, fill=azul!8, rounded corners=4pt,
                     minimum width=3.8cm, minimum height=0.6cm,
                     inner sep=4pt, align=center},
        arr/.style={-Stealth, thick, color=verde}]

        \node[fase] (A) {\icono{comment-dots} El niño habla con OLIBOT};
        \node[fase, below=0.3cm of A] (B) {\icono{brain} OLIBOT elige la actividad};
        \node[fase, below=0.3cm of B] (C) {\icono{pencil-alt} El niño traza / responde};
        \node[fase, below=0.3cm of C] (D) {\icono{star} OLIBOT celebra y adapta};
        \node[fase, below=0.3cm of D] (E) {\icono{arrow-right} Avanza al siguiente reto};

        \draw[arr] (A) -- (B);
        \draw[arr] (B) -- (C);
        \draw[arr] (C) -- (D);
        \draw[arr] (D) -- (E);
        \draw[arr] (E.east) -- ++(0.3,0) |- (A.east);

        \node[font=\tiny, color=gray, anchor=west]
          at ([xshift=0.4cm, yshift=0.1cm]A.east) {loop adaptativo};
      \end{tikzpicture}
    \end{column}
  \end{columns}
\end{frame}

%% ─── 6. ACTIVIDADES ─────────────────────────────────────────
\begin{frame}{¿Qué hace el niño durante la sesión?}

  \begin{columns}[T]
    \begin{column}{.32\textwidth}
      \begin{block}{\centering\icono{pencil-alt}\; Trazado de letras}
        \small\centering
        El niño sigue los puntos guía en la pantalla táctil.\\[4pt]
        Tres niveles de ayuda: con guía completa, con puntos clave, y libre.\\[4pt]
        \textbf{Demostración automática} antes de que el niño intente.
      \end{block}
    \end{column}
    \begin{column}{.32\textwidth}
      \begin{exampleblock}{\centering\icono{microphone}\; Interacción por voz}
        \small\centering
        El niño escucha las instrucciones de OLIBOT y responde hablando.\\[4pt]
        No requiere saber leer ni escribir para interactuar.\\[4pt]
        Compatible con la voz \textbf{en español de España}.
      \end{exampleblock}
    \end{column}
    \begin{column}{.32\textwidth}
      \begin{alertblock}{\centering\icono{palette}\; Colorear y dibujar}
        \small\centering
        Como actividad de descanso cognitivo o premio, el programa ofrece fichas de colorear digitales.\\[4pt]
        Mantiene la \textbf{motivación} durante toda la sesión.
      \end{alertblock}
    \end{column}
  \end{columns}

  \vspace{10pt}
  \begin{center}
    \begin{tikzpicture}
      \node[fill=verde!15, rounded corners=6pt, inner sep=8pt,
            text width=12cm, align=center, font=\small] {
        \icono{info-circle}\;
        \textbf{Sesiones de 15--20 minutos.}
        El docente permanece en el aula en todo momento.
        El programa \textbf{no sustituye} a la maestra: la \textbf{complementa}.
      };
    \end{tikzpicture}
  \end{center}
\end{frame}

%% ============================================================
\section{Enfoque pedagógico}
%% ============================================================

%% ─── 7. PEDAGOGÍA ───────────────────────────────────────────
\begin{frame}{Principios pedagógicos de OLIBOT}

  \begin{columns}[T]
    \begin{column}{.48\textwidth}
      \begin{block}{\icono{layer-group}\; Scaffolding (Vygotsky)}
        \footnotesize
        Ayudas graduadas ante dificultades:
        \begin{enumerate}
          \item \textit{«¿Quieres intentarlo otra vez?»}
          \item \textit{«Fíjate en el punto de inicio…»}
          \item Demostración completa del trazo
        \end{enumerate}
        Nunca da la respuesta sin que el niño lo intente.
      \end{block}
      \begin{block}{\icono{user-check}\; Perfil individual}
        \footnotesize
        Cada niño tiene su \textbf{nombre y progreso} guardado. El sistema recuerda qué letras necesita reforzar.
      \end{block}
    \end{column}
    \begin{column}{.48\textwidth}
      \begin{exampleblock}{\icono{shield-alt}\; Seguridad para menores}
        \footnotesize
        \begin{itemize}
          \item Solo responde sobre las actividades (contenido filtrado).
          \item Sin acceso a internet ni grabación de audio.
          \item Lenguaje siempre positivo.
        \end{itemize}
      \end{exampleblock}
      \begin{exampleblock}{\icono{star}\; Motivación intrínseca}
        \footnotesize
        \begin{itemize}
          \item Celebraciones y refuerzo positivo.
          \item Avatar animado que «siente» el progreso.
          \item Insignias al superar bloques de contenido.
          \item Progresión visible para el niño.
        \end{itemize}
      \end{exampleblock}
    \end{column}
  \end{columns}
\end{frame}

%% ============================================================
\section{Plan de evaluación}
%% ============================================================

%% ─── 8. QUÉ QUEREMOS MEDIR ──────────────────────────────────
\begin{frame}{¿Qué queremos evaluar?}

  \begin{center}
  \begin{tikzpicture}[font=\small,
    metrica/.style={draw=azul, fill=azul!8, rounded corners=6pt,
                    inner sep=8pt, text width=4.0cm, align=left,
                    drop shadow={opacity=0.15}},
    titulo/.style={font=\bfseries\color{azul}}]

    % Fila 1
    \node[metrica] (M1) {
      {\titulo\icono{clock}\; Atención sostenida}\\[3pt]
      Tiempo de interacción activa con OLIBOT medido por sesión y por alumno.};

    \node[metrica, right=0.4cm of M1] (M2) {
      {\titulo\icono{chart-bar}\; Rendimiento comparado}\\[3pt]
      Actividades superadas \textbf{con} OLIBOT frente a fichas en papel equivalentes (grupo control).};

    \node[metrica, right=0.4cm of M2] (M3) {
      {\titulo\icono{hand-paper}\; Intervención del adulto}\\[3pt]
      N.º de veces que el docente o investigadora intervienen para ayudar al niño.};

    % Fila 2
    \node[metrica, below=0.4cm of M1] (M4) {
      {\titulo\icono{smile}\; Satisfacción del usuario}\\[3pt]
      Encuesta post-sesión con \textbf{caras} (\faIcon{smile}/\faIcon{frown}) y observación de actitud.};

    \node[metrica, right=0.4cm of M4] (M5) {
      {\titulo\icono{graduation-cap}\; Velocidad de progresión}\\[3pt]
      Letras/sílabas nuevas superadas por sesión vs.\ línea base de la tutora.};

    \node[metrica, right=0.4cm of M5] (M6) {
      {\titulo\icono{redo}\; Engagement voluntario}\\[3pt]
      N.º de veces que el niño pide «otra vez» o intenta un trazo de nuevo sin pedírselo.};

  \end{tikzpicture}
  \end{center}

  \vspace{4pt}
  \begin{center}
    \small\textcolor{gray}{%
    Los datos se recogen de forma anónima, sin identificación personal del alumno.}
  \end{center}
\end{frame}

%% ─── 9. DISEÑO EVALUACIÓN ───────────────────────────────────
\begin{frame}{Diseño de la evaluación}

  \begin{columns}[T]
    \begin{column}{.50\textwidth}
      \begin{block}{\icono{calendar-alt}\; Cronograma propuesto}
        \footnotesize
        \begin{tabular}{@{}ll@{}}
          \toprule
          \textbf{Fase} & \textbf{Duración} \\
          \midrule
          Sesión de presentación al equipo  & 1 sesión \\
          Evaluación con alumnos            & 3--5 sesiones \\
          \quad — Grupo usando OLIBOT        & 15--20 min/sesión \\
          \quad — Grupo control (papel)      & misma duración \\
          Recogida de feedback docente       & 1 sesión \\
          Entrega de informe al centro       & 1 semana después \\
          \bottomrule
        \end{tabular}
      \end{block}
      \begin{alertblock}{\icono{exclamation-circle}\; Importante}
        \footnotesize
        El docente \textbf{siempre estará presente} y puede interrumpir la sesión en cualquier momento.
      \end{alertblock}
    \end{column}
    \begin{column}{.46\textwidth}
      \begin{exampleblock}{\icono{users}\; Muestra}
        \footnotesize
        \begin{itemize}
          \item Alumnos de \textbf{4 o 5 años}.
          \item Grupo reducido de \textbf{4--8 niños}.
          \item División aleatoria: OLIBOT vs.\ control.
        \end{itemize}
      \end{exampleblock}
      \begin{exampleblock}{\icono{tools}\; Material necesario}
        \footnotesize
        \begin{itemize}
          \item \textbf{1 tablet} con Chrome o Safari.
          \item \textbf{WiFi} o punto de acceso propio.
          \item Espacio tranquilo en el aula.
        \end{itemize}
      \end{exampleblock}
    \end{column}
  \end{columns}
\end{frame}

%% ============================================================
\section{Garantías y compromisos}
%% ============================================================

%% ─── 10. GARANTÍAS ──────────────────────────────────────────
\begin{frame}{Lo que garantizamos}

  \begin{columns}[T]
    \begin{column}{.48\textwidth}
      \begin{block}{\icono{user-shield}\; Privacidad y protección de datos}
        \small
        \begin{itemize}
          \item \textbf{Sin nombres ni datos personales}: los perfiles usan apodos o iniciales.
          \item \textbf{Sin grabaciones de audio}: el reconocimiento de voz se procesa localmente.
          \item Cumplimiento de la \textbf{LOPDGDD} y normativa de protección de menores.
          \item Los datos de uso se anonimizarán antes de cualquier análisis.
        \end{itemize}
      \end{block}
      \vspace{4pt}
      \begin{block}{\icono{university}\; Contexto académico}
        \small
        Este trabajo es parte de un \textbf{TFM} de la Universidad Carlos III de Madrid. Los resultados se usarán exclusivamente con fines de investigación académica, \textbf{sin ningún uso comercial}.
      \end{block}
    \end{column}
    \begin{column}{.48\textwidth}
      \begin{exampleblock}{\icono{hands-helping}\; Lo que recibe el centro}
        \small
        \begin{itemize}
          \item[\tick] Acceso gratuito a OLIBOT durante la evaluación.
          \item[\tick] \textbf{Informe} del progreso por alumno (anónimo).
          \item[\tick] Informe docente con sugerencias pedagógicas.
          \item[\tick] Sin compromiso de uso posterior.
          \item[\tick] Uso gratuito continuado si lo desean.
        \end{itemize}
      \end{exampleblock}
      \begin{center}
        \begin{tikzpicture}
          \node[fill=verde, rounded corners=6pt, inner sep=8pt,
                text=white, font=\bfseries] {%
            \icono{envelope}\; Firmamos carta de compromiso si lo desean};
        \end{tikzpicture}
      \end{center}
    \end{column}
  \end{columns}
\end{frame}

%% ─── 11. LO QUE PEDIMOS ─────────────────────────────────────
\begin{frame}{Lo que necesitamos de vosotros}

  \begin{center}
  \begin{tikzpicture}[font=\small,
    item/.style={draw=verde, fill=verde!8, rounded corners=6pt,
                 inner sep=7pt, text width=2.4cm, align=center,
                 minimum height=2.3cm, drop shadow={opacity=0.12}}]

    \node[item] (P1) {
      \Large\icono{calendar-check}\\[4pt]
      \small\textbf{3--5 sesiones}\\[1pt]
      \footnotesize 15--20\,min por sesión};

    \node[item, right=0.4cm of P1] (P2) {
      \Large\icono{chalkboard-teacher}\\[4pt]
      \small\textbf{Tutora\\presente}\\[1pt]
      \footnotesize observa y garantiza el bienestar};

    \node[item, right=0.4cm of P2] (P3) {
      \Large\icono{wifi}\\[4pt]
      \small\textbf{Conexión WiFi}\\[1pt]
      \footnotesize o punto de acceso propio};

    \node[item, right=0.4cm of P3] (P4) {
      \Large\icono{file-signature}\\[4pt]
      \small\textbf{Autorización}\\[1pt]
      \footnotesize familias y dirección (modelo estándar)};

    \node[item, right=0.4cm of P4] (P5) {
      \Large\icono{comment-alt}\\[4pt]
      \small\textbf{Feedback}\\[1pt]
      \footnotesize cuestionario breve al docente};

  \end{tikzpicture}
  \end{center}

  \vspace{0.2cm}
  \begin{center}
    \begin{tikzpicture}
      \node[fill=azul!10, draw=azul, rounded corners=5pt, inner sep=6pt,
            text width=12.5cm, align=center, font=\footnotesize] {
        \textbf{En ningún caso interrumpiremos el horario lectivo regular.}
        Las sesiones se realizarán en los momentos que la tutora considere más oportunos
        (rincones, trabajo en grupos, tiempo libre guiado…)
      };
    \end{tikzpicture}
  \end{center}
\end{frame}

%% ============================================================
\section{Conclusión y demostración}
%% ============================================================

%% ─── 12. RESUMEN BENEFICIOS ─────────────────────────────────
\begin{frame}{¿Por qué vale la pena probar OLIBOT?}

  \begin{columns}[T]
    \begin{column}{.54\textwidth}
      \begin{exampleblock}{\icono{check-double}\; Beneficios directos para el aula}
        \normalsize
        \begin{itemize}
          \item El docente \textbf{gana tiempo} para atender a quienes más lo necesitan mientras OLIBOT acompaña a otro grupo.
          \item Cada niño avanza \textbf{a su ritmo} sin presión ni comparación.
          \item Refuerza exactamente \textbf{lo que se trabaja en clase}.
          \item Genera datos objetivos que el docente puede usar para su \textbf{evaluación continua}.
        \end{itemize}
      \end{exampleblock}
    \end{column}
    \begin{column}{.42\textwidth}
      \begin{tikzpicture}[font=\small]
        \node[draw=verde, fill=verde!8, rounded corners=6pt,
              inner sep=10pt, text width=4.5cm, align=center] (cita) {
          \large\itshape\color{azul}
          «La tecnología bien usada no reemplaza al docente:\\
          \textbf{multiplica su impacto}»\\[6pt]
          \footnotesize\color{gray}
          Adaptado de Salman Khan,\\
          \textit{The One World Schoolhouse}, 2012
        };
      \end{tikzpicture}
      \vspace{8pt}
      \begin{block}{\icono{eye}\; Veámoslo en vivo}
        \small
        A continuación realizamos una \textbf{demostración} de exactamente lo que vería un alumno en su primera sesión.\\[4pt]
        La demo dura \textbf{5 minutos}.
      \end{block}
    \end{column}
  \end{columns}
\end{frame}

%% ─── 13. DIAPOSITIVA DE DEMO ────────────────────────────────
\begin{frame}[plain]
  \begin{tikzpicture}[remember picture, overlay]
    \fill[verde] (current page.south west) rectangle (current page.north east);

    \node[white, font=\bfseries\fontsize{48}{52}\selectfont]
      at ([yshift=1.2cm]current page.center) {DEMO EN VIVO};

    \node[white, font=\Large]
      at (current page.center) {\icono{tablet-alt}\;\; Se inicia OLIBOT…};

    \node[fill=white, rounded corners=8pt, inner sep=12pt,
          font=\normalsize\color{verde}\bfseries,
          text width=7cm, align=center]
      at ([yshift=-2.2cm]current page.center)
      {\faIcon{info-circle}\; Recordatorio para la demo:\\[4pt]
       \footnotesize\color{azul}
       1. Seleccionar nombre y edad del niño\\
       2. Mostrar la demo automática del trazo\\
       3. Demostrar interacción por voz\\
       4. Mostrar la adaptación de nivel};
  \end{tikzpicture}
\end{frame}

%% ─── 14. CIERRE ─────────────────────────────────────────────
\begin{frame}{Próximos pasos}

  \begin{columns}[T]
    \begin{column}{.50\textwidth}
      \begin{block}{\icono{list-ol}\; Si decide participar…}
        \normalsize
        \begin{enumerate}
          \item \textbf{Hoy}: comentar la propuesta con el equipo docente.
          \item \textbf{Esta semana}: firmar carta de colaboración + autorización familiar.
          \item \textbf{Próximas semanas}: acordar fechas de sesiones.
          \item \textbf{Tras las sesiones}: recibirá el informe completo.
        \end{enumerate}
      \end{block}
    \end{column}
    \begin{column}{.46\textwidth}
      \begin{exampleblock}{\icono{address-card}\; Contacto}
        \normalsize
        \textbf{María Emilia Núñez Guerrero}\\[4pt]
        \icono{envelope}\; \texttt{menemezgu@alumnos.uc3m.es}\\[4pt]
        \icono{university}\; Departamento de Informática\\
        Universidad Carlos III de Madrid\\[8pt]
        \small\textcolor{gray}{Supervisado por:\\
        Dr.\ Javier Ignacio Carbó Rubiera}
      \end{exampleblock}
    \end{column}
  \end{columns}

  \vspace{14pt}
  \begin{center}
    \begin{tikzpicture}
      \node[fill=verde, rounded corners=8pt, inner sep=14pt,
            text=white, font=\bfseries\Large] {
        \faIcon{heart}\;\; ¡Muchas gracias por su tiempo y confianza!};
    \end{tikzpicture}
  \end{center}
\end{frame}

%% ============================================================
%% APÉNDICE (no se muestra en demo, disponible si hay preguntas)
%% ============================================================
\appendix

\begin{frame}{Preguntas frecuentes — Privacidad}
  \small
  \begin{block}{¿Se graba la voz de los niños?}
    No. El reconocimiento de voz usa la API del propio navegador del dispositivo.
    Los datos de audio \textbf{no se almacenan} en ningún servidor.
  \end{block}
  \vspace{4pt}
  \begin{block}{¿Qué datos se recogen exactamente?}
    Únicamente: número de intentos por actividad, nivel de pistas solicitadas y tiempo de sesión.
    \textbf{Ningún dato permite identificar individualmente al alumno} (se usan apodos o iniciales elegidos por el docente).
  \end{block}
  \vspace{4pt}
  \begin{block}{¿Necesita el niño usuario o contraseña?}
    No. El docente selecciona el perfil del niño antes de que se siente ante la tablet.
    El niño no necesita ningún acceso digital previo.
  \end{block}
\end{frame}

\begin{frame}{Alineación curricular detallada}
  \small
  \begin{columns}[T]
    \begin{column}{.48\textwidth}
      \begin{exampleblock}{3 años — Pre-grafomotricidad}
        \begin{itemize}
          \item Trazos verticales, horizontales y curvos
          \item Conciencia fonológica oral
          \item Reconocimiento de vocales A, E
        \end{itemize}
      \end{exampleblock}
      \vspace{4pt}
      \begin{exampleblock}{4 años — Vocales y primeras consonantes}
        \begin{itemize}
          \item Todas las vocales (A, E, I, O, U)
          \item Consonantes m, l, s, p, t
          \item Números del 1 al 5
        \end{itemize}
      \end{exampleblock}
    \end{column}
    \begin{column}{.48\textwidth}
      \begin{exampleblock}{5 años — Sílabas y palabras}
        \begin{itemize}
          \item Consonantes d, n, r, f, b, c, g
          \item Sílabas directas (ma, me, mi, mo, mu)
          \item Primeras palabras (mamá, papá, sol)
          \item Números del 1 al 10
        \end{itemize}
      \end{exampleblock}
      \vspace{4pt}
      \begin{block}{Referencia normativa}
        \footnotesize
        Decreto 36/2022, de 9 de mayo (CM).\\
        Área: «Comunicación y Representación de la Realidad».\\
        Criterios de evaluación 4.1, 4.2, 4.3.
      \end{block}
    \end{column}
  \end{columns}
\end{frame}

\end{document}
